\section{Markov 演化、半群与跳过程}\label{sec:06-03}

\subsection{Markov 核、Chapman--Kolmogorov 与转移半群}\label{subsec:6-3-1}
\index{Markov 核}\index{Chapman--Kolmogorov 方程}\index{转移半群}

知道过程的当前状态，并不总能预测下一步。若
\[
 X_{n+1}=X_n+X_{n-1}+\xi_{n+1},
\]
其中 $\xi_{n+1}$ 是新的独立噪声，那么给定 $X_n$ 仍不足以确定下一步的条件分布；至少还要记住
$X_{n-1}$。当然可以把状态扩成 $(X_{n-1},X_n)$，但这也说明了 Markov 性的真正内容：所选的
“当前状态”必须已经包含预测未来所需的全部过去信息。正则条件分布只保证条件分布可以写成核，
却不会替我们证明这一充分性。

先看三个核，它们依次把确定演化、离散噪声与扩散放进同一格式。

\begin{example}[确定性动力系统]\label{ex:6-3-1-1}
设 $(E,\mathcal E)$ 为可测空间，$T:E\to E$ 可测。定义
\[
 K(x,A)=\delta_{T(x)}(A)=\1_A(T(x)).
\]
固定 $x$ 时这是概率测度；固定 $A$ 时它关于 $x$ 可测。对有界可测 $f$，
\[
 Kf(x)=\int_E f(y)K(x,\dd y)=f(T(x)).
\]
因此确定性轨道 $x,T(x),T^2(x),\ldots$ 是 Markov 演化的退化情形，且
$K^nf=f\circ T^n$。
\end{example}

\begin{example}[加性噪声]\label{ex:6-3-1-2}
令 $E=\R^d$，$F:E\to E$ 为 Borel 映射，$\xi_1,\xi_2,\ldots$ 独立同分布，共同分布为
$\nu$，并与 $X_0$ 独立。递推式
\[
 X_{n+1}=F(X_n)+\xi_{n+1}
\]
给出
\[
 K(x,A)=\nu(A-F(x)),\qquad
 Kf(x)=\int_{\R^d}f(F(x)+z)\nu(\dd z).
\]
可测性先对 $f=\1_A$ 由 $(x,z)\mapsto F(x)+z$ 的联合可测性和 Tonelli 得到；简单函数与
单调收敛再给一般非负 $f$。独立性与取出已知因子说明
$\E[f(X_{n+1})\mid\mathcal F_n]=Kf(X_n)$，其中
$\mathcal F_n=\sigma(X_0,\xi_1,\ldots,\xi_n)$。
\end{example}

\begin{example}[Gaussian 热核]\label{ex:6-3-1-3}
沿用定义~\ref{def:4-4-2-2} 的热核 $p_t$，并令
\[
 q_t(x)=p_{t/2}(x)=(2\pi t)^{-d/2}e^{-|x|^2/(2t)}.
\]
它是 $N(0,tI_d)$ 的密度。由
$p_s*p_t=p_{s+t}$ 得 $q_s*q_t=q_{s+t}$。于是
\[
 P_tf(x)=\int_{\R^d}f(x+y)q_t(y)\dd y
\]
满足 $P_sP_t=P_{s+t}$，保持正函数和常数，并且
$\|P_tf\|_\infty\leq\|f\|_\infty$。按照本小节稍后在
评注~\ref{proposition:6-3-1-4} 中给出的生成元定义，这里的时间规范使生成元在光滑函数上为
$\tfrac12\Delta$；第 4 章的 $p_t$ 规范则给生成元 $\Delta$。
\end{example}

\begin{definition}[Markov 核]\label{def:6-3-1-1}
从可测空间 $(E,\mathcal E)$ 到 $(F,\mathcal F)$ 的 \emph{Markov 核}是函数
$K:E\times\mathcal F\to[0,1]$，满足：
\begin{enumerate}
 \item 对每个 $x\in E$，$A\mapsto K(x,A)$ 是 $(F,\mathcal F)$ 上的概率测度；
 \item 对每个 $A\in\mathcal F$，$x\mapsto K(x,A)$ 是 $\mathcal E$-可测函数。
\end{enumerate}
若 $f:F\to\R$ 非负或有界可测，定义
\[
 Kf(x)=\int_F f(y)K(x,\dd y).
\]
若 $\mu$ 是 $E$ 上的概率，定义
\[
 \mu K(A)=\int_EK(x,A)\mu(\dd x).
\]
前一个运算把函数从 $F$ 拉回 $E$，后一个运算把测度从 $E$ 推向 $F$；两者方向相反。
\end{definition}

\begin{definition}[Markov 性]\label{def:6-3-1-2}
设 $(X_n)$ 对滤过 $(\mathcal F_n)$ 适应并取值于 $(E,\mathcal E)$。若存在核
$K:E\leadsto E$，使每个有界可测 $f:E\to\R$ 都满足
\begin{equation}\label{eq:6-3-1-markov-property}
 \E[f(X_{n+1})\mid\mathcal F_n]=Kf(X_n)\qquad\text{a.s.},
\end{equation}
则称 $X$ 是相对于该滤过、转移核为 $K$ 的时间齐次 Markov 链。时间非齐次情形以
单步核 $K_n$ 表示从时刻 $n$ 到 $n+1$ 的转移，即把右端改为 $K_nf(X_n)$。
\end{definition}

多步核不能仅凭某个给定初始分布下的条件期望版本，在未到达状态上逐点识别；因此多步核取为单步核的复合。

若已经指定一个从 $x$ 出发、满足 $\Prob_x(X_0=x)=1$ 的 Markov 分布，就把相应的概率与期望
记为 $\Prob_x$、$\E_x$；初始分布为 $\mu$ 时相应记为 $\Prob_\mu$、$\E_\mu$。这套下标记号只记录
所选初始分布，不自行断言该分布的存在或唯一性。

\begin{remark}[条件分布与标准 Borel 假设]\label{remark:6-3-1-rcd}
若 $E$ 是标准 Borel 空间，定理~\ref{theorem:5-2-2-1} 可为 $X_{n+1}$ 给定 $X_n$ 选择
正则条件分布；这说明可以选出一个以当前状态为参数的可测核。式~\eqref{eq:6-3-1-markov-property} 还要求给定整个
$\mathcal F_n$ 后没有额外变化，这一要求不是正则条件分布存在性的推论。若扩大滤过并泄露未来噪声，
原来的 Markov 等式也可能失效。
若 $E$ 不是标准 Borel 空间，条件期望等式仍可作为 Markov 性的定义，但把逐个集合得到的条件概率
同时选成一个关于状态联合可测的核，可能做不到。例~\ref{ex:6-3-1-2} 的显式公式
$K(x,A)=\nu(A-F(x))$ 同时绕过了这两层选择问题。
\end{remark}

核可以逐步拼接。下面的证明把“对集合可测”和“对函数积分”两层都写出来。

\begin{proposition}[核的复合]\label{proposition:6-3-1-1}
若 $K:E\leadsto F$、$L:F\leadsto G$ 是 Markov 核，则
\[
 (KL)(x,A)=\int_F L(y,A)K(x,\dd y),\qquad A\in\mathcal G,
\]
定义一个核 $KL:E\leadsto G$。对每个非负或有界可测 $f:G\to\R$，
\[
 (KL)f=K(Lf).
\]
\end{proposition}

\begin{proof}
固定 $A\in\mathcal G$。函数 $y\mapsto L(y,A)$ 可测且非负；核积分的简单函数逼近说明
$x\mapsto(KL)(x,A)$ 可测。固定 $x$，若 $(A_j)$ 两两不交，则 Tonelli 定理给
\begin{align*}
 (KL)\left(x,\bigcup_jA_j\right)
 &=\int_F\sum_jL(y,A_j)K(x,\dd y)\\
 &=\sum_j\int_FL(y,A_j)K(x,\dd y)
 =\sum_j(KL)(x,A_j).
\end{align*}
此外 $(KL)(x,G)=1$，故 $A\mapsto(KL)(x,A)$ 为概率。

对 $f=\1_A$，函数恒等式就是定义。线性性把它推广到非负简单函数，单调收敛再推广到非负
可测函数；有界实值函数对正负部应用该结论，复值函数再对实部、虚部应用。
\end{proof}

\begin{theorem}[Chapman--Kolmogorov 方程]\label{theorem:6-3-1-2}
设时间非齐次 Markov 链的单步核为 $(K_n)_{n\ge0}$。置
\[
 K_{m,m}=I,\qquad K_{m,n}=K_mK_{m+1}\cdots K_{n-1}\quad(m<n).
\]
则对每个有界可测 $f$，
\[
 \E[f(X_n)\mid\mathcal F_m]=K_{m,n}f(X_m)\quad\text{a.s.}\qquad(m\le n),
\]
并且这些逐点定义的核满足
\[
 K_{m,n}=K_{m,k}K_{k,n}\qquad(m\leq k\leq n).
\]
时间齐次离散链的 $n$ 步核为 $K^n$；若 $X_0$ 的分布是 $\mu$，则
\[
 \Law(X_n)=\mu K^n,
 \qquad
 \E_\mu f(X_n)=\int_EK^nf(x)\mu(\dd x)
\]
对每个有界可测 $f$ 成立。
\end{theorem}

\begin{proof}
先核验核复合的结合律。对任意三个可复合核 $K,L,R$ 及任意有界可测 $f$，
命题~\ref{proposition:6-3-1-1} 给出
\[
 ((KL)R)f=(KL)(Rf)=K(L(Rf))=K((LR)f)=(K(LR))f.
\]
取 $f=\1_A$，便有 $((KL)R)(x,A)=(K(LR))(x,A)$ 对每个 $x,A$ 成立。因此上式中有限个
单步核的复合无须加括号，并且按 $m\le k\le n$ 分组立即得到
$K_{m,n}=K_{m,k}K_{k,n}$；这是逐点核恒等式，不受某个初始分布的零质量状态影响。

再证明条件期望公式。$n=m$ 时公式是适应性给出的
$\E[f(X_m)\mid\mathcal F_m]=f(X_m)$；$n=m+1$ 时就是 Markov 性。若公式对
$(m+1,n)$ 成立，则塔式性质及单步 Markov 等式给
\begin{align*}
 \E[f(X_n)\mid\mathcal F_m]
 &=\E\bigl[K_{m+1,n}f(X_{m+1})\mid\mathcal F_m\bigr]\\
 &=K_m(K_{m+1,n}f)(X_m)=K_{m,n}f(X_m).
\end{align*}
对 $n-m$ 归纳即得一般公式。

时间齐次时对 $n$ 归纳得 $K_{0,n}=K^n$。最后
\[
 \Prob_\mu(X_n\in A)
 =\int_EK^n(x,A)\mu(\dd x)=(\mu K^n)(A),
\]
函数公式由简单函数与单调收敛得到。
\end{proof}

上面的证明把多步预测压成终点函数。若还要观察中间路径，核仍可递归积分。

\begin{proposition}[Markov 条件期望的路径版本]\label{proposition:6-3-1-3}
设 $X$ 是转移核为 $K$ 的时间齐次 Markov 链。对 $m\geq1$ 和有界可测
$F:E^m\to\R$，存在有界可测函数 $G_F:E\to\R$，使
\[
 \E[F(X_{n+1},\ldots,X_{n+m})\mid\mathcal F_n]=G_F(X_n)\quad\text{a.s.}
\]
具体地，$G_F(x)$ 由从最末坐标向前逐次对 $K$ 积分得到。
\end{proposition}

\begin{proof}
先把“逐次积分”写成无歧义的公式。对有界可测
$F:E^m\to\R$，定义
\begin{equation}\label{eq:6-3-1-path-kernel}
 G_F(x_0)=
 \int_E\!\cdots\!\int_E
 F(x_1,\ldots,x_m)
 K(x_{m-1},\dd x_m)\cdots K(x_0,\dd x_1).
\end{equation}
该公式可从最后一个变量向前归纳理解。具体地，参数核积分的可测性与
命题~\ref{proposition:6-3-1-1} 的证明相同：对示性函数成立，再经简单函数和单调收敛
推广。逐层使用这一事实，可知式~\eqref{eq:6-3-1-path-kernel} 的每一层积分都是有界可测
函数，特别地 $G_F$ 有界可测。

先取矩形乘积函数
$F(x_1,\ldots,x_m)=\prod_{j=1}^mf_j(x_j)$，其中每个 $f_j$ 有界可测。令
\[
 g_m=f_m,\qquad g_j(x)=f_j(x)K g_{j+1}(x)\quad(j=m-1,\ldots,1).
\]
此时式~\eqref{eq:6-3-1-path-kernel} 给出 $G_F=Kg_1$。从 $X_{n+m}$ 开始，反复使用塔式性质和
式~\eqref{eq:6-3-1-markov-property}，得到
\[
 \E\!\left[\prod_{j=1}^mf_j(X_{n+j})\middle|\mathcal F_n\right]
 =(Kg_1)(X_n)=G_F(X_n).
\]

令 $\mathcal H$ 为使所述条件等式以式~\eqref{eq:6-3-1-path-kernel} 中的 $G_F$ 成立的有界可测
$F$ 的集合。条件期望与核积分的线性性说明 $\mathcal H$ 是向量空间。若
$0\le F_j\uparrow F$ 且整列一致有界，则条件单调收敛和式
\eqref{eq:6-3-1-path-kernel} 中的逐层单调收敛同时给出 $F\in\mathcal H$。矩形示性函数已属于
$\mathcal H$，函数单调类定理遂给出全部有界可测 $F$。
\end{proof}

确定时刻的 Markov 性还可以停在一个随机时刻，但离散性在证明中承担了关键的可数分片。

\begin{theorem}[离散时间强 Markov 性]\label{theorem:6-3-1-strong-markov}
设 $X$ 是相对于 $(\mathcal F_n)$ 的时间齐次 Markov 链，$\tau<\infty$ a.s. 是同一滤过的停时。
则对 $m\geq0$ 与有界可测 $f$，
\[
 \E[f(X_{\tau+m})\mid\mathcal F_\tau]=K^mf(X_\tau)\quad\text{a.s.}
\]
更一般地，停时后的任意有界有限柱函数由命题~\ref{proposition:6-3-1-3} 的递归核从
$X_\tau$ 起计算。
\end{theorem}

\begin{proof}
记 $Y=K^mf(X_\tau)$。它对 $\mathcal F_\tau$ 可测：对 Borel 集 $B$ 与整数 $n$，
\[
 \{Y\in B\}\cap\{\tau\leq n\}
 =\bigcup_{k=0}^n\{K^mf(X_k)\in B\}\cap\{\tau=k\}\in\mathcal F_n.
\]
任取 $A\in\mathcal F_\tau$。由停止信息的定义，
$A\cap\{\tau=n\}\in\mathcal F_n$。确定时刻 Markov 性给
\begin{align*}
 \E[\1_Af(X_{\tau+m})]
 &=\sum_{n=0}^\infty
   \E[\1_{A\cap\{\tau=n\}}f(X_{n+m})]\\
 &=\sum_{n=0}^\infty
   \E[\1_{A\cap\{\tau=n\}}K^mf(X_n)]
 =\E[\1_AY].
\end{align*}
级数可交换是因为 $f$ 有界而各片两两不交。这个积分恒等式正是条件期望的定义。

有限未来柱函数先对矩形函数逐层使用同一计算，再以函数单调类推广。证明依赖的是可数值停时的分片；
连续时间强 Markov 性还需要路径连续性和滤过右连续性，不能从这里直接类推。
\end{proof}

\begin{definition}[转移半群]\label{def:6-3-1-3}
设 $(K_t)_{t\geq0}$ 是 $E$ 上的 Markov 核族，满足
\[
 K_0=I,\qquad K_{s+t}=K_sK_t\quad(s,t\geq0).
\]
若适应过程 $X=(X_t)_{t\geq0}$ 对每个有界可测 $f$ 及 $0\leq s\leq t$ 满足
\[
 \E[f(X_t)\mid\mathcal F_s]=K_{t-s}f(X_s)\quad\text{a.s.},
\]
则称 $X$ 是相对于 $(\mathcal F_t)$、转移核族为 $(K_t)$ 的连续时间齐次 Markov 过程。
定义算子 $P_tf=K_tf$。若对每个 $x$ 还指定了满足
$\Prob_x(X_0=x)=1$ 的过程分布，且
$K_t(x,A)=\Prob_x(X_t\in A)$，则同一算子也写成
$P_tf(x)=\E_xf(X_t)$。核的 Chapman--Kolmogorov 恒等式给出
$P_0=I$ 与 $P_{s+t}=P_sP_t$。每个 $P_t$ 保持非负函数和常数 $1$，并满足
$\|P_tf\|_\infty\leq\|f\|_\infty$；这样的算子族称为转移半群或 Markov 半群。
\end{definition}

Poisson 过程、一致化构造与 Brownian motion 的转移半群都可由独立增量或显式核验证这一恒等式。
连续时间的半群律是额外结构，不是离散时间定理~\ref{theorem:6-3-1-2} 的形式推论。

\begin{remark}[Feller 性与生成元的边界]\label{proposition:6-3-1-4}
若 $E$ 为局部紧 Hausdorff 空间，且 $P_tC_0(E)\subset C_0(E)$，则称 $(P_t)$ 具有 Feller 性。
若进一步对每个 $f\in C_0(E)$ 有 $\|P_tf-f\|_\infty\to0$，它便是第 4 章意义下的收缩
$C_0$ 半群，可在
\[
 D(L)=\left\{f:\lim_{t\downarrow0}\frac{P_tf-f}{t}
 \text{ 在 }C_0(E)\text{ 中存在}\right\}
\]
定义生成元。半群律本身既不推出 Feller 性，也不推出强连续性。有限链和一致有界跳率链提供了可直接验证的例子；
一般 Feller 过程的构造与生成元刻画则需要更强的假设。
\end{remark}

\begin{figure}[htbp]
\centering
\begin{tikzpicture}[node distance=13mm and 18mm]
 \node[book strong node] (x) {状态 $x$};
 \node[book node,right=of x] (y) {中间状态 $y$};
 \node[book strong node,right=of y] (z) {状态 $z$};
 \draw[book accent arrow] (x)--node[above,book label]{$K_s(x,\dd y)$}(y);
 \draw[book accent arrow] (y)--node[above,book label]{$K_t(y,\dd z)$}(z);
 \draw[book dashed arrow,bend right=28] (x) to
   node[below,book label]{$K_{s+t}(x,\dd z)=\int K_t(y,\dd z)K_s(x,\dd y)$}(z);
 \node[book warm node,below=16mm of y] (f) {函数方向\\$P_{s+t}f=P_s(P_tf)$};
 \draw[book arrow] (z)--(f);
 \draw[book arrow] (f)--(x);
\end{tikzpicture}
\caption{核沿时间正向复合；函数从终点向起点反向积分，同一恒等式写成半群律。}
\label{fig:6-3-1-kernel-semigroup}
\end{figure}

这个结构有三条容易混淆的边界。第一，$Kf$ 和 $\mu K$ 的方向相反。第二，Markov 性相对于给定滤过
表述；加入未来信息可能破坏它。第三，半群律只组织有限维转移分布，不保证样本路径连续，甚至不保证
$t\mapsto P_tf$ 强连续。

\begin{exercise}
\begin{enumerate}
 \item 逐项验证加性噪声核 $K(x,A)=\nu(A-F(x))$ 的两个核条件。
 \item 证明 $K^nf(x)=\E_xf(X_n)$，并写出 $K^2(x,A)$ 的二重积分。
 \item 对热核 $q_t$ 直接完成平方，重新证明 $q_s*q_t=q_{s+t}$。
 \item 构造一个 Markov 链及一个泄露 $X_{n+1}$ 的扩大滤过，使链不再满足相同的 Markov 条件等式。
\end{enumerate}
\end{exercise}

核可以有不变分布，链也可能从任意初态逐渐忘掉起点；这两句话并不等价。下一小节把不变、
分布收敛和时间平均分开处理。

\subsection{不变分布、可逆性与 Markov 链遍历}\label{subsec:6-3-2}
\index{不变分布}\index{平稳 Markov 链}\index{可逆性}\index{遍历定理}

若 $X_0$ 的分布为 $\mu$，则 $X_n$ 的分布是 $\mu P^n$。一个自然问题是寻找
$\pi P=\pi$；然而即使找到了这样的 $\pi$，仍有两件事没有回答：从给定初态出发的
$P^n(i,\cdot)$ 是否趋于 $\pi$，以及单条路径的时间平均是否趋于 $\pi$-平均。周期链会否定第一句
而保留第二句，多个闭沟通类则会同时破坏唯一性和全局收敛。

\begin{definition}[不变分布与平稳链]\label{def:6-3-2-1}
若状态空间上的概率 $\pi$ 满足 $\pi P=\pi$，则称 $\pi$ 是转移核 $P$ 的\emph{不变分布}。
若 $X_0\sim\pi$，则每个 $X_n\sim\pi$；更进一步，对任意
$0\leq n_1<\cdots<n_k$ 与 $r\geq0$，Markov 核的迭代公式给
\[
 \Law(X_{n_1},\ldots,X_{n_k})
 =\Law(X_{n_1+r},\ldots,X_{n_k+r}).
\]
为了看清最后一步，对任意可测集 $A_1,\ldots,A_k$，路径核公式给出
\begin{align*}
 &\Prob(X_{n_1+r}\in A_1,\ldots,X_{n_k+r}\in A_k)\\
 &\quad=\int_E\!\cdots\!\int_E\pi(\dd x_0)
   P^{n_1+r}(x_0,\dd x_1)\1_{A_1}(x_1)
   \prod_{j=2}^k P^{n_j-n_{j-1}}(x_{j-1},\dd x_j)\1_{A_j}(x_j)\\
 &\quad=\int_E\!\cdots\!\int_E\pi(\dd x_1)\1_{A_1}(x_1)
   \prod_{j=2}^k P^{n_j-n_{j-1}}(x_{j-1},\dd x_j)\1_{A_j}(x_j),
\end{align*}
其中第二个等号使用 $\pi P^{n_1+r}=\pi$；对未平移的时刻组同样使用
$\pi P^{n_1}=\pi$ 得到同一积分。矩形是乘积 $\sigma$-代数的生成 $\pi$-系统，因而上式确定整个
联合分布。这样的链称为平稳链。平稳只表示分布对时间平移不变，不表示各时刻相互独立。
\end{definition}

\begin{definition}[可逆性]\label{def:6-3-2-2}
若 $\pi$ 是状态空间上的概率，并且 $E\times E$ 上有测度恒等
\begin{equation}\label{eq:6-3-2-detailed-balance}
 \pi(\dd x)P(x,\dd y)=\pi(\dd y)P(y,\dd x),
\end{equation}
则称 $P$ 相对于 $\pi$ \emph{可逆}。在有限或可数状态空间上，这等价于细致平衡（detailed balance）关系
\[
 \pi(i)P(i,j)=\pi(j)P(j,i)\qquad(i,j\in E).
\]
对式~\eqref{eq:6-3-2-detailed-balance} 的两边关于 $x$ 积分，得到
$\pi P=\pi$；所以可逆性蕴含不变性，反向一般不成立。
\end{definition}

\begin{definition}[不可约与周期]\label{def:6-3-2-3}
有限状态链称为\emph{不可约}，若任意 $i,j$ 之间都存在 $n\geq0$ 使 $P^n(i,j)>0$。
状态 $i$ 的周期为
\[
 d(i)=\gcd\{n\geq1:P^n(i,i)>0\}.
\]
不可约链中所有状态的周期相同；共同周期为 $1$ 时称链非周期。
\end{definition}

上述“共同周期”需要由不可约性推出。任取状态 $i,j$，选择
$r,s\ge0$ 使 $P^r(i,j)>0$ 且 $P^s(j,i)>0$。若 $n\ge1$ 且 $P^n(j,j)>0$，则
Chapman--Kolmogorov 方程给出
\[
 P^{r+s}(i,i)\ge P^r(i,j)P^s(j,i)>0,
\]
以及
\[
 P^{r+n+s}(i,i)
 \ge P^r(i,j)P^n(j,j)P^s(j,i)>0.
\]
因而 $d(i)$ 同时整除 $r+s$ 和 $r+n+s$，也就整除 $n$。这对每个满足
$P^n(j,j)>0$ 的 $n$ 成立，故 $d(i)\mid d(j)$。交换 $i,j$ 得
$d(j)\mid d(i)$，从而 $d(i)=d(j)$。

\begin{example}[二状态链]\label{ex:6-3-2-1}
令
\[
 P=\begin{pmatrix}1-a&a\\ b&1-b\end{pmatrix},
 \qquad 0\leq a,b\leq1,\quad a+b>0.
\]
解 $\pi P=\pi$ 与 $\pi_0+\pi_1=1$，得到
\[
 \pi=\left(\frac b{a+b},\frac a{a+b}\right).
\]
向量 $(1,-1)$ 是特征值 $1-a-b$ 的左特征向量；任意两个概率行向量之差都与
$(1,-1)$ 成比例，因而
\[
 \mu P^n-\pi=(1-a-b)^n(\mu-\pi).
\]
当 $a,b>0$ 时链不可约。若 $a=b=1$，链在两个状态间确定性交替，周期为 $2$，所以
$P^n(0,\cdot)$ 不收敛；但其 \Cesaro{} 平均趋于 $(1/2,1/2)$。其余不可约情形满足
$|1-a-b|<1$，上式给出几何收敛率。
\end{example}

\begin{example}[有限无向图上的随机游走]\label{ex:6-3-2-2}
设 $G=(V,E_G)$ 是有限连通无向图，$d(x)$ 为顶点度。简单随机游走的转移概率为
$P(x,y)=d(x)^{-1}$（若 $x\sim y$），否则为零。令
\[
 \pi(x)=\frac{d(x)}{\sum_{z\in V}d(z)}.
\]
当 $x\sim y$ 时
\[
 \pi(x)P(x,y)=\frac1{\sum_zd(z)}=\pi(y)P(y,x),
\]
其余情形两边同为零，所以链可逆。若图为二分图，每一步都从一个部跳到另一个部，周期为 $2$；
把转移改为 $(I+P)/2$，每个状态都有正的自环概率，周期便降为 $1$，不变分布保持不变。
\end{example}

\begin{example}[多个闭类]\label{ex:6-3-2-3}
在状态集 $\{0,1\}$ 上取 $P=I$。两个状态都是吸收态，每个
$\theta\delta_0+(1-\theta)\delta_1$ 都是不变分布；从 $i$ 出发的分布恒为 $\delta_i$。
所以“不变分布存在”既不保证唯一，也不保证极限与初态无关。缺失的正是不可约性。
\end{example}

\begin{example}[无限状态中不可约仍不够]\label{ex:6-3-2-4}
在 $\Z$ 上的简单对称随机游走不可约。若 $\pi$ 是不变概率，则
\[
 \pi(k)=\frac{\pi(k-1)+\pi(k+1)}2,
\]
故差分 $\pi(k+1)-\pi(k)$ 为常数。一个在整条 $\Z$ 上非负的仿射序列只能是常数，而非零常数
不可求和，所以不存在不变概率。将一步改为以 $1/2$ 留在原地、以 $1/4$ 各向左右移动，只把上式
两边同时加入 $\pi(k)/2$，结论不变；此链已经非周期。计数测度虽然可逆，却有无限总质量。
这说明有限状态结论不能只凭“不可约、非周期”移植到一般状态空间；还必须另行证明足以控制
回返概率与回返时间的条件。
\end{example}

\begin{figure}[htbp]
\centering
\begin{tikzpicture}[node distance=12mm and 18mm]
  \node[book strong node] (a0) {$0$};
  \node[book strong node,right=of a0] (a1) {$1$};
  \draw[book accent arrow,bend left=20] (a0) to node[above,book label]{1} (a1);
  \draw[book accent arrow,bend left=20] (a1) to node[below,book label]{1} (a0);
  \node[book label,above=7mm of a0,xshift=9mm] {闭类 $A={0,1}$：周期 $2$};

  \node[book warm node,right=34mm of a1] (b) {$2$};
  \draw[book accent arrow,loop above] (b) to node[above,book label]{1} (b);
  \node[book label,above=7mm of b] {闭类 $B={2}$：周期 $1$};

  \node[book node,below=18mm of a1,xshift=19mm] (tr) {瞬时态 $3$};
  \draw[book dashed arrow] (tr) to node[left,book label]{正概率} (a1);
  \draw[book dashed arrow] (tr) to node[right,book label]{正概率} (b);
\end{tikzpicture}
\caption{两个闭沟通类使不变分布不唯一；左侧的二循环又显示，即使限制在一个不可约闭类内，
周期仍可使单时刻分布来回振荡。定理~\ref{theorem:6-3-2-1}--\ref{theorem:6-3-2-3}
分别给出不可约情形下的不变性、单时刻收敛与时间平均。}
\label{fig:6-3-2-classes-period}
\end{figure}

有限状态不可约性提供一个很实用的几何尾估计。固定状态 $i_*$。对每个 $x$ 选择一条从 $x$
到 $i_*$ 的正概率有限路径。状态有限，因而存在 $m<\infty$ 与 $\varepsilon>0$，使从任意 $x$
出发在 $m$ 步以内命中 $i_*$ 的概率至少为 $\varepsilon$。更明确地，设为 $x$ 选取的路径长度为
$m_x$、路径概率乘积为 $p_x>0$，则可取
\[
 m=\max_xm_x,\qquad \varepsilon=\min_xp_x>0.
\]
若 $\tau_{i_*}=\inf\{n\geq0:X_n=i_*\}$，则 Markov 性在确定时刻 $km$ 给出
\begin{align*}
 \Prob_x(\tau_{i_*}>(k+1)m)
 &=\E_x\!\left[
   \1_{\{\tau_{i_*}>km\}}
   \Prob_{X_{km}}(\tau_{i_*}>m)
   \right]\\
 &\le(1-\varepsilon)\Prob_x(\tau_{i_*}>km).
\end{align*}
从 $k=0$ 归纳便得
\begin{equation}\label{eq:6-3-2-hitting-geometric-tail}
 \sup_x\Prob_x(\tau_{i_*}>km)\leq(1-\varepsilon)^k.
\end{equation}
特别地，所有命中时间和从 $i_*$ 出发的首次正回返时间
$\tau_{i_*}^+=\inf\{n\geq1:X_n=i_*\}$ 都有有限期望。
由于从 $i_*$ 出发的 $\tau_{i_*}$ 等于零，式~\eqref{eq:6-3-2-hitting-geometric-tail}
本身不直接控制正回返时间。由第一步后的 Markov 性，
若 $X_1=y$，则 $\tau_{i_*}^+$ 的剩余部分与从 $y$ 出发命中 $i_*$ 的时间同分布，故
\[
 \E_{i_*}\tau_{i_*}^+
 =1+\sum_{y\in E}P(i_*,y)\E_y\tau_{i_*}<\infty.
\]
这里每一项由式~\eqref{eq:6-3-2-hitting-geometric-tail} 可积，且状态集有限。

\begin{theorem}[有限不可约链的不变分布]\label{theorem:6-3-2-1}
有限状态不可约 Markov 链存在唯一的不变分布 $\pi$，并且 $\pi(j)>0$ 对每个状态 $j$ 成立。
固定 $i\in E$ 时，它由一次回返周期的占用时间给出：
\begin{equation}\label{eq:6-3-2-cycle-invariant-law}
 \pi(j)=
 \frac{\E_i\sum_{k=0}^{\tau_i^+-1}\1_{\{X_k=j\}}}
      {\E_i\tau_i^+}.
\end{equation}
\end{theorem}

\begin{proof}
式~\eqref{eq:6-3-2-hitting-geometric-tail} 说明分母有限且为正。记
\[
 u(j)=\E_i\sum_{k=0}^{\tau_i^+-1}\1_{\{X_k=j\}},
 \qquad m_i=\E_i\tau_i^+.
\]
由于 $\sum_ju(j)=m_i$，式~\eqref{eq:6-3-2-cycle-invariant-law} 定义一个概率。对任意 $\ell$，
把随机上限写成非负指标函数后，条件期望和 Tonelli 定理给
\begin{align*}
 \sum_ju(j)P(j,\ell)
&=\sum_{k=0}^\infty
  \E_i\!\left[\1_{\{k<\tau_i^+\}}P(X_k,\ell)\right]\\
&=\sum_{k=0}^\infty
  \E_i\!\left[\1_{\{k<\tau_i^+\}}
   \E_i[\1_{\{X_{k+1}=\ell\}}\mid\mathcal F_k]\right]\\
 &=\E_i\sum_{k=1}^{\tau_i^+}\1_{\{X_k=\ell\}}.
\end{align*}
这里 $\{k<\tau_i^+\}\in\mathcal F_k$，所以第三行取出已知因子的步骤合法。
最后一个和与 $u(\ell)$ 的差为
$\1_{\{X_{\tau_i^+}=\ell\}}-\1_{\{X_0=\ell\}}=0$，因为周期首尾都等于 $i$。
故 $uP=u$，从而 $\pi P=\pi$。从 $i$ 到 $j$ 可选择一条内部不再经过 $i$ 的简单正概率路径，
所以一个回返周期以正概率访问 $j$，即 $u(j)>0$。

只剩唯一性。若 $\rho P=\rho$，定义反向转移矩阵
\[
 \widehat P(\ell,j)=\frac{\pi(j)P(j,\ell)}{\pi(\ell)}.
\]
不变性说明每一行和为一；不可约性在反向边上保持。令 $h(j)=\rho(j)/\pi(j)$，则
\[
 h(\ell)=\sum_j\widehat P(\ell,j)h(j).
\]
取 $h$ 的最大点 $\ell_0$。上式是若干不超过 $h(\ell_0)$ 的数的凸组合，所以每个满足
$\widehat P(\ell_0,j)>0$ 的 $j$ 都取同一最大值。沿反向链的正概率路径传播，不可约性迫使
$h$ 在全部状态上为常数。由 $\sum_j\rho(j)=\sum_j\pi(j)=1$，该常数为 $1$，故 $\rho=\pi$。
\end{proof}

非周期性可以转成所有状态共享的一块正质量。

\begin{lemma}[有限不可约非周期矩阵的共同正幂]\label{lemma:6-3-2-primitive}
若有限状态转移矩阵 $P$ 不可约且非周期，则存在 $m\geq1$，使
$P^m(i,j)>0$ 对所有 $i,j$ 成立。
\end{lemma}

\begin{proof}
固定状态 $r$。集合 $A=\{n\geq1:P^n(r,r)>0\}$ 对加法封闭，且其最大公因数为 $1$。
可从 $A$ 中选有限多个 $a_1,\ldots,a_q$，其最大公因数仍为 $1$。初等整数论说明由这些
$a_k$ 的非负整数线性组合组成的半群包含所有充分大的整数。具体地，其余生成元的剩余类在
$\Z/a_1\Z$ 中生成整个有限群；有限群中一个生成元剩余类的负元是该生成元剩余类的某个非负整数倍，
所以每个剩余类 $c$ 都可选到一个由 $a_2,\ldots,a_q$ 的非负组合表示的整数 $b_c\ge0$。若
$n\equiv c\pmod{a_1}$ 且 $n\ge\max_cb_c$，则
$n=b_c+\ell a_1$，其中 $\ell\ge0$。这就证明所有充分大的整数都属于该半群。

对每个 $i,j$，由不可约性选择正概率路径 $i\to r$ 和 $r\to j$，长度分别记为
$\alpha_i,\beta_j$。当 $n-\alpha_i-\beta_j$ 足够大时，可在 $r$ 处插入若干回路，使总长度恰为
$n$，因而 $P^n(i,j)>0$。状态对只有有限多个，取一个同时充分大的 $m$ 即得结论。
\end{proof}

\begin{theorem}[有限遍历收敛]\label{theorem:6-3-2-2}
若有限状态 Markov 链不可约且非周期，$\pi$ 为其不变分布，则对每个初态 $i$，
\[
 d_{\mathrm{TV}}(P^n(i,\cdot),\pi)\longrightarrow0.
\]
收敛事实上具有几何速度。
\end{theorem}

\begin{proof}
取引理~\ref{lemma:6-3-2-primitive} 中的 $m$，令
\[
 a(j)=\min_iP^m(i,j),\qquad
 \eta=\sum_ja(j)>0,
 \qquad \nu(j)=a(j)/\eta.
\]
若 $\eta=1$，所有 $P^m(i,\cdot)$ 都等于 $\nu$，结论直接成立。以下设 $0<\eta<1$，并定义
\[
 R(i,j)=\frac{P^m(i,j)-\eta\nu(j)}{1-\eta}.
\]
$R$ 是转移矩阵，且 $P^m(i,\cdot)=\eta\nu+(1-\eta)R(i,\cdot)$。因此对任意概率
$\mu,\rho$，Markov 核的全变差收缩性给
\[
 d_{\mathrm{TV}}(\mu P^m,\rho P^m)
 =(1-\eta)d_{\mathrm{TV}}(\mu R,\rho R)
 \leq(1-\eta)d_{\mathrm{TV}}(\mu,\rho).
\]
这里最后一个不等式可直接从
$\sup_A|\int R(x,A)(\mu-\rho)(\dd x)|$ 及有符号测度的正负分解得到。
迭代后
\[
 d_{\mathrm{TV}}(\delta_iP^{qm},\pi)
 \leq(1-\eta)^q d_{\mathrm{TV}}(\delta_i,\pi),
\]
因为 $\pi P^{qm}=\pi$。若 $n=qm+r$、$0\leq r<m$，再作用 $P^r$ 只会缩小全变差，故结论成立。
\end{proof}

单时刻收敛用到了非周期性；时间平均不需要它。回返时刻把一条相关路径切成独立周期。

\begin{theorem}[时间平均遍历定理]\label{theorem:6-3-2-3}
若 $(X_n)$ 是有限状态不可约 Markov 链，则对任意函数 $f:E\to\R$ 及任意初始分布，
\[
 \frac1n\sum_{k=0}^{n-1}f(X_k)\longrightarrow\sum_{x\in E}f(x)\pi(x)
 \qquad\text{a.s.}
\]
其中 $\pi$ 是定理~\ref{theorem:6-3-2-1} 的唯一不变分布。
\end{theorem}

\begin{proof}
固定状态 $i$。从任意初态出发，式~\eqref{eq:6-3-2-hitting-geometric-tail} 保证首次命中 $i$
的时刻 $\tau_0$ 有限 a.s.。递归定义
\[
 \tau_{r+1}=\inf\{n>\tau_r:X_n=i\},\qquad
 L_r=\tau_{r+1}-\tau_r,
 \qquad
 Y_r=\sum_{k=\tau_r}^{\tau_{r+1}-1}f(X_k).
\]
离散强 Markov 性说明 $(L_r,Y_r)_{r\geq0}$ 独立同分布，其分布等于从 $i$ 出发的一个回返
周期。这一点可以逐项验证：对 $\ell\ge1$，事件
\[
 \{L_r=\ell,\ (X_{\tau_r},\ldots,X_{\tau_r+\ell})\in A\}
\]
是停时后有限柱事件；强 Markov 性说明它在给定 $\mathcal F_{\tau_r}$ 后的概率只等于从
$i$ 出发的相应概率。对 $\ell$ 取可数并，再对 $r$ 归纳，便得整个周期块两两独立且同分布；
$Y_r$ 是该有限块的可测函数，所以结论传递到 $(L_r,Y_r)$。几何尾界给
$\E L_0<\infty$；状态有限使 $f$ 有界，故
$|Y_0|\leq\|f\|_\infty L_0$ 也可积。强大数律给
\[
 \frac{\sum_{r=0}^{q-1}Y_r}{\sum_{r=0}^{q-1}L_r}\longrightarrow
 \frac{\E_i\sum_{k=0}^{\tau_i^+-1}f(X_k)}{\E_i\tau_i^+}
 =\sum_xf(x)\pi(x).
\]

需要控制首尾不完整周期。首块长度 $\tau_0$ 是有限随机数，除以 $n$ 后趋零。对尾块，
可积性和尾和公式给每个 $\delta>0$
\[
 \sum_{r=1}^\infty\Prob(L_r>\delta r)
 =\sum_{r=1}^\infty\Prob(L_0>\delta r)
 \le \frac1\delta\int_0^\infty\Prob(L_0>u)\dd u+1
 =\frac{\E L_0}{\delta}+1<\infty,
\]
所以 Borel--Cantelli 说明 $L_r/r\to0$ a.s.。若
$N(n)=\max\{r:\tau_r\leq n\}$，则周期长的强大数律给
$\tau_r/r\to\E L_0$。几何尾界还说明每个回返时刻有限，所以 $N(n)\to\infty$。置
$r=N(n)$，由 $\tau_r\le n<\tau_{r+1}$ 得
\[
 \frac{r}{\tau_r}\ge\frac r n,
 \qquad
 \frac{r+1}{n}>\frac{r+1}{\tau_{r+1}}.
\]
两端分别趋于 $1/\E L_0$，而 $(r+1)/n-r/n=1/n\to0$；夹逼给
$N(n)/n\to1/\E L_0$。因而
$L_{N(n)}/n=(L_{N(n)}/N(n))(N(n)/n)\to0$。尾块奖励的绝对值不超过
$\|f\|_\infty L_{N(n)}$。更明确地，对充分大的 $n$，
\[
 \sum_{k=0}^{n-1}f(X_k)
 =\sum_{k=0}^{\tau_0-1}f(X_k)
  +\sum_{r=0}^{N(n)-1}Y_r
  +\sum_{k=\tau_{N(n)}}^{n-1}f(X_k).
\]
首项除以 $n$ 趋零，末项绝对值除以 $n$ 至多
$\|f\|_\infty L_{N(n)}/n\to0$。又
\[
 \frac{\tau_{N(n)}-\tau_0}{n}
 =1-\frac{n-\tau_{N(n)}+\tau_0}{n}\longrightarrow1,
\]
因为 $0\le n-\tau_{N(n)}<L_{N(n)}$。把完整周期的奖励和写成
\[
 \frac1n\sum_{r=0}^{N(n)-1}Y_r
 =\frac{\sum_{r=0}^{N(n)-1}Y_r}{\sum_{r=0}^{N(n)-1}L_r}
  \frac{\tau_{N(n)}-\tau_0}{n},
\]
前一因子已由强大数律确定，故全时间平均收敛到 $\sum_xf(x)\pi(x)$。
\end{proof}

\begin{proposition}[周期链的 \Cesaro{} 分布收敛]\label{proposition:6-3-2-5}
有限不可约链无论周期为何，都有
\[
 \frac1n\sum_{k=0}^{n-1}P^k(i,\cdot)\longrightarrow\pi
 \quad\text{于全变差}.
\]
\end{proposition}

\begin{proof}
对每个状态 $j$，把遍历定理用于 $f=\1_{\{j\}}$，再用有界收敛定理取期望，得到
\[
 \frac1n\sum_{k=0}^{n-1}P^k(i,j)\longrightarrow\pi(j).
\]
状态集有限，逐坐标收敛可在有限和中取极限；由
$d_{\mathrm{TV}}(\mu,\nu)=\tfrac12\sum_j|\mu(j)-\nu(j)|$ 即得全变差收敛。
二状态确定性交替链说明原来的 $P^n(i,\cdot)$ 可以不收敛。
\end{proof}

\begin{proposition}[可逆算子的 $L^2$ 结构]\label{proposition:6-3-2-4}
若 $P$ 相对于 $\pi$ 可逆，则它在实或复 Hilbert 空间 $L^2(\pi)$ 上自伴且为收缩：
\[
 \langle f,Pg\rangle_\pi=\langle Pf,g\rangle_\pi,
 \qquad \|Pf\|_{L^2(\pi)}\leq\|f\|_{L^2(\pi)}.
\]
\end{proposition}

\begin{proof}
采用 $\langle f,g\rangle_\pi=\int\overline f g\,\dd\pi$ 的复内积约定。对有界 $f,g$，
细致平衡恒等式与 Fubini 定理给出
\begin{align*}
 \langle f,Pg\rangle_\pi
 &=\int\!\int \overline{f(x)}g(y)P(x,\dd y)\pi(\dd x)\\
 &=\int\!\int \overline{f(x)}g(y)P(y,\dd x)\pi(\dd y)
 =\langle Pf,g\rangle_\pi.
\end{align*}
截断后令截断高度趋于无穷，可推广到 $L^2$。另一方面，Jensen 不等式和不变性给
\[
 \|Pf\|_2^2
 \leq\int P(|f|^2)(x)\pi(\dd x)
 =\int|f|^2\dd\pi.
\]
\end{proof}

自伴性为谱方法打开了入口，却不能单独消除周期。二状态确定性交替链在零均值子空间上的特征值为
$-1$；只控制“第二大特征值”而忽略负端，不能推出 $P^n$ 收敛。惰性化把它改成 $0$，或者必须控制
零均值子空间上的绝对谱半径。一般谱隙理论留待后文，这里不以它重新证明有限遍历结论。

\begin{remark}[MCMC 中的三种长时间问题]
在有限状态 MCMC 中，通常先设计核 $P$ 使目标分布 $\pi$ 满足细致平衡，由
定义~\ref{def:6-3-2-2} 得到 $\pi P=\pi$。这只解决“目标分布是否不变”。若还有不可约性与
非周期性，定理~\ref{theorem:6-3-2-2} 才给出从固定初态出发的全变差收敛，这是讨论预热期（burn-in）的
分布依据。无论周期如何，定理~\ref{theorem:6-3-2-3} 仍保证单条轨道上的样本平均收敛。
因此，“不变”、“预热后近似目标分布”和“时间平均可用”是三个需要不同结论支撑的陈述。
\end{remark}

\begin{figure}[htbp]
\centering
\begin{tikzpicture}[node distance=9mm and 12mm]
 \node[book strong node] (i0) {回返 $i$\ $\tau_0$};
 \node[book node,right=of i0] (b1) {周期 1\\长度 $L_0$\\奖励 $Y_0$};
 \node[book strong node,right=of b1] (i1) {回返 $i$\ $\tau_1$};
 \node[book node,right=of i1] (b2) {周期 2\\长度 $L_1$\\奖励 $Y_1$};
 \node[book strong node,right=of b2] (i2) {回返 $i$\ $\tau_2$};
 \draw[book accent arrow] (i0)--(b1);
 \draw[book accent arrow] (b1)--(i1);
 \draw[book accent arrow] (i1)--(b2);
 \draw[book accent arrow] (b2)--(i2);
 \node[book warm node,below=15mm of i1] (law) {强 Markov 性 $\Longrightarrow (L_r,Y_r)$ i.i.d.};
 \draw[book dashed arrow] (b1)--(law);
 \draw[book dashed arrow] (b2)--(law);
\end{tikzpicture}
\caption{回返把相关的 Markov 路径切成独立同分布的再生块；时间平均成为奖励和与长度和的比。}
\label{fig:6-3-2-regeneration}
\end{figure}

\begin{exercise}
\begin{enumerate}
 \item 求二状态链的 $P^n$，并分别验证周期情形的单时刻失败和 \Cesaro{} 收敛。
 \item 证明细致平衡蕴含不变性，并给出一个不变但不可逆的有限链。
 \item 对有限无向连通图核验度分布；判断三角形、偶环和加入自环后的周期。
 \item 在定理~\ref{theorem:6-3-2-3} 的证明中，逐步证明
 $N(n)/n\to1/\E_i\tau_i^+$。
\end{enumerate}
\end{exercise}

有限链的长期结构由回返组织。连续时间最基本的回返模型从随机等待开始：若等待时间没有记忆，
计数过程就会产生独立平稳增量。

\subsection{Poisson 过程、跳链与生成元}\label{subsec:6-3-3}
\index{Poisson 过程}\index{指数等待时间}\index{补偿鞅}\index{$Q$-矩阵}

\begin{example}[稀有 Bernoulli 极限]\label{ex:6-3-3-1}
把长度为 $t$ 的时间段切成 $n$ 个小格，并假设每格以概率 $\lambda t/n$ 发生一次事件，格间相互
独立。总事件数 $B_n$ 服从 $\operatorname{Bin}(n,\lambda t/n)$，所以对固定 $k$，
\[
 \Prob(B_n=k)
 =\binom nk\left(\frac{\lambda t}{n}\right)^k
  \left(1-\frac{\lambda t}{n}\right)^{n-k}
 \longrightarrow e^{-\lambda t}\frac{(\lambda t)^k}{k!}.
\]
这个计算给出了 Poisson 边缘，却没有说明不同时间段的计数怎样耦合。过程的独立增量、等待时间和
路径阶梯结构还必须另外建立。
\end{example}

\begin{definition}[Poisson 过程]\label{def:6-3-3-1}
速率 $\lambda>0$ 的 \emph{Poisson 过程}是过程 $(N_t)_{t\geq0}$，满足：
\begin{enumerate}
 \item $N_0=0$ a.s.，路径右连续有左极限、非减、取非负整数值且每次跳跃大小为 $1$；
 \item 不相交时间区间上的增量独立，且增量平稳；
 \item 对 $0\leq s<t$，
 \[
   N_t-N_s\sim\operatorname{Poisson}(\lambda(t-s)).
 \]
\end{enumerate}
\end{definition}

若先从这些有限维分布在典范乘积空间上扩张，随后构造的阶梯路径过程只保证与之有相同的过程分布；
除非两者已置于同一概率空间并逐时刻比较，不能把“同分布的正则实现”称为修正。

\begin{definition}[到达时刻与等待时间]\label{def:6-3-3-2}
对单位跳计数过程，令
\[
 T_0=0,\qquad T_n=\inf\{t\geq0:N_t\geq n\},\qquad
 W_n=T_n-T_{n-1}.
\]
$T_n$ 是第 $n$ 次到达时刻，$W_n$ 是相邻到达间隔。若所有 $T_n$ 有限且趋于无穷，则
\[
 N_t=\max\{n\geq0:T_n\leq t\}.
\]
\end{definition}

等待时间的无记忆性是 Markov 性在单个时钟上的原型。

\begin{proposition}[指数分布的无记忆性]\label{proposition:6-3-3-1}
若 $W\sim\operatorname{Exp}(\lambda)$，即
$\Prob(W>t)=e^{-\lambda t}$（$t\geq0$），则
\[
 \Prob(W>s+t\mid W>s)=\Prob(W>t).
\]
反过来，设 $W$ 取值于扩展半轴 $[0,\infty]$，生存函数 $G(t)=\Prob(W>t)$ 连续，并满足
$G(s+t)=G(s)G(t)$。若 $W$ 既不 a.s. 等于 $0$，也不 a.s. 等于 $+\infty$，则存在
$\lambda>0$ 使 $G(t)=e^{-\lambda t}$。
\end{proposition}

\begin{proof}
指数分布的正向计算为
\[
 \frac{\Prob(W>s+t)}{\Prob(W>s)}
 =\frac{e^{-\lambda(s+t)}}{e^{-\lambda s}}=e^{-\lambda t}.
\]
反向中，乘法式在 $s=t=0$ 时给 $G(0)=G(0)^2$。若 $G(0)=0$，则
$\Prob(W>0)=0$，即 $W=0$ a.s.，这是被排除的退化情形；故 $G(0)=1$。非退化性与连续性还
排除某个有限时刻取零：若 $G(t_0)=0$，则
$G(t_0/n)^n=0$，所以 $G(t_0/n)=0$ 对每个 $n$，与 $G(0)=1$ 处连续矛盾。因此
$H(t)=-\log G(t)$ 定义良好、连续、非降并满足 $H(s+t)=H(s)+H(t)$。对有理
$q\geq0$，加法性给 $H(q)=qH(1)$；连续性把它推广到全部 $t\geq0$。令
$\lambda=H(1)$。若 $\lambda=0$，则 $G\equiv1$，对应 $W=+\infty$ a.s.，已被排除；故
$\lambda>0$。
\end{proof}

\begin{theorem}[由等待时间构造 Poisson 过程]\label{theorem:6-3-3-2}
设 $W_1,W_2,\ldots$ 独立同分布为 $\operatorname{Exp}(\lambda)$，令
\[
 T_n=\sum_{k=1}^nW_k,\qquad
 N_t=\max\{n\geq0:T_n\leq t\}.
\]
则 $T_n\to\infty$ a.s.；在这个概率一事件上，上式对每个有限 $t$ 都定义一个有限整数，
并在其补集上统一取零路径。这样得到的阶梯过程是速率
$\lambda$ 的 Poisson 过程。
\end{theorem}

\begin{proof}
事件 $A_n=\{W_n\geq1\}$ 独立且概率为 $e^{-\lambda}>0$，所以
$\sum_n\Prob(A_n)=\infty$。Borel--Cantelli II 说明 $A_n$ 发生无穷多次，故
$T_n\to\infty$ a.s.。把例外零集上的整条路径统一定义为零以后，每个 $N_t$ 都处处有定义，
而且后续全部路径性质可在同一个概率一事件上同时核验。

卷积归纳给 $T_n$ 的 Gamma 密度
\begin{equation}\label{eq:6-3-3-gamma-density}
 f_{T_n}(u)=\frac{\lambda^n u^{n-1}}{(n-1)!}e^{-\lambda u}\1_{(0,\infty)}(u).
\end{equation}
确切地说，$n=1$ 是指数密度；若公式对 $n$ 成立，则
\begin{align*}
 f_{T_{n+1}}(u)
 &=\int_0^u\frac{\lambda^nv^{n-1}}{(n-1)!}e^{-\lambda v}
     \lambda e^{-\lambda(u-v)}\dd v\\
 &=\frac{\lambda^{n+1}u^n}{n!}e^{-\lambda u}.
\end{align*}
于是对 $n\geq0$，
\begin{align*}
 \Prob(N_t=n)
 &=\Prob(T_n\leq t<T_{n+1})\\
 &=\int_0^tf_{T_n}(u)\Prob(W_{n+1}>t-u)\dd u
 =e^{-\lambda t}\frac{(\lambda t)^n}{n!},
\end{align*}
其中 $n=0$ 直接由 $\Prob(W_1>t)$ 计算。

还要证明增量独立。固定确定时刻 $s$，记
$\mathcal F_s^N=\sigma(N_u:0\le u\le s)$。在 $\{N_s=n\}$ 上，$s$ 以前的整条计数路径由
$(T_1,\ldots,T_n)$ 决定；反过来，每个 $T_j\wedge s$ 都是该路径的可测函数。因此
$\mathcal F_s^N$ 在 $\{N_s=n\}$ 上的迹由事件
\[
 \{N_s=n\}\cap\{(T_1,\ldots,T_n)\in C\},
 \qquad C\in\Borel(\R^n),
\]
生成（$n=0$ 时只剩 $\{N_s=0\}$）。

在 $\{N_s=n\}$ 上置剩余等待时间 $R_s=T_{n+1}-s$。设 $r\ge0$，
$C\in\Borel(\R^n)$，并取 Borel 集 $A_2,\ldots,A_m\subset(0,\infty)$。对
$(W_1,\ldots,W_n)$ 条件化，并使用后续等待时间的独立性，得到
\begin{align*}
&\Prob\bigl((T_1,\ldots,T_n)\in C,\ N_s=n,\ R_s>r,
             \ W_{n+j}\in A_j\ (2\le j\le m)\bigr)\\
&=\E\!\left[
 \1_{\{(T_1,\ldots,T_n)\in C,\ T_n\le s\}}
 \Prob(W_{n+1}>s-T_n+r\mid W_1,\ldots,W_n)
 \right]
 \prod_{j=2}^m\Prob(W_1\in A_j)\\
&=e^{-\lambda r}
 \E\!\left[
 \1_{\{(T_1,\ldots,T_n)\in C,\ T_n\le s\}}
 e^{-\lambda(s-T_n)}\right]
 \prod_{j=2}^m\Prob(W_1\in A_j)\\
&=e^{-\lambda r}\prod_{j=2}^m\Prob(W_1\in A_j)\,
 \Prob\bigl((T_1,\ldots,T_n)\in C,\ N_s=n\bigr).
\end{align*}
第一行中的 $N_s=n$ 等价于 $T_n\le s<T_{n+1}$；最后一行使用
$\Prob(W_{n+1}>s-T_n\mid W_1,\ldots,W_n)=e^{-\lambda(s-T_n)}$。
对矩形 $\prod_{j=2}^mA_j$ 先用有限交，再用 $\pi$--$\lambda$ 定理，以上分解说明：条件于
$\mathcal F_s^N$，$R_s,W_{N_s+2},W_{N_s+3},\ldots$ 相互独立、都服从
$\operatorname{Exp}(\lambda)$，而且其联合条件分布不依赖过去。于是从 $s$ 起的到达过程由一列
新的 i.i.d. 指数等待时间构造，并与 $\mathcal F_s^N$ 独立。已经证明的边缘计算
应用于这列新时钟，得到
$N_{s+t}-N_s\sim\operatorname{Poisson}(\lambda t)$，并由递归得到任意有限个不相交区间的
增量独立。阶梯路径的右连续、左极限、非减和单位跳性质由定义直接读出。
\end{proof}

反方向不能只由第一到达时间的尾公式猜测全部等待时间；必须同时读出不同到达之间的联合结构。

\begin{proposition}[Poisson 增量反推出指数等待]\label{proposition:6-3-3-arrivals}
对定义~\ref{def:6-3-3-1} 的 Poisson 过程，等待时间
$W_1,W_2,\ldots$ 独立同分布为 $\operatorname{Exp}(\lambda)$。更精确地，
$(T_1,\ldots,T_n)$ 在区域 $0<t_1<\cdots<t_n$ 上具有联合密度
\begin{equation}\label{eq:6-3-3-arrival-joint-density}
 \lambda^ne^{-\lambda t_n}\1_{\{0<t_1<\cdots<t_n\}}.
\end{equation}
\end{proposition}

\begin{proof}
先固定 $t>0$ 和分割 $0=s_0<s_1<\cdots<s_m=t$。独立增量给出，在
$\sum_jn_j=n$ 时
\begin{align*}
 &\Prob\bigl(N_{s_j}-N_{s_{j-1}}=n_j,\ 1\leq j\leq m\mid N_t=n\bigr)\\
 &\qquad=
 \frac{n!}{\prod_jn_j!}\prod_{j=1}^m
 \left(\frac{s_j-s_{j-1}}{t}\right)^{n_j}.
\end{align*}
这正是 $n$ 个独立 $\operatorname{Unif}(0,t)$ 点落入各小区间的多项分布。半开区间组成的有限
分割生成 $[0,t]^n$ 的 Borel $\sigma$-代数，因此条件于 $N_t=n$，到达点的无序配置与
$n$ 个独立均匀点相同；排序后，
\[
 (T_1,\ldots,T_n)\mid\{N_t=n\}
\]
在 $0<t_1<\cdots<t_n<t$ 上具有密度 $n!/t^n$。

为严格得到无条件密度，固定 $T>0$。条件于 $N_T=m\ge n$ 时，全部 $m$ 个到达点是 $m$ 个
$\operatorname{Unif}(0,T)$ 点的次序统计量。把后 $m-n$ 个有序坐标在
$t_n<u_{n+1}<\cdots<u_m<T$ 上积分，得到前 $n$ 个到达时刻的条件密度
\[
 \frac{m!}{(m-n)!T^m}(T-t_n)^{m-n}
 \1_{\{0<t_1<\cdots<t_n<T\}}.
\]
乘以 $\Prob(N_T=m)=e^{-\lambda T}(\lambda T)^m/m!$ 并对 $m\ge n$ 求和，Tonelli 定理给出
\begin{align*}
 &\sum_{m=n}^\infty
 e^{-\lambda T}\frac{(\lambda T)^m}{m!}
 \frac{m!}{(m-n)!T^m}(T-t_n)^{m-n}\\
 &\qquad=e^{-\lambda T}\lambda^n
 \sum_{r=0}^\infty\frac{(\lambda(T-t_n))^r}{r!}
 =\lambda^ne^{-\lambda t_n}.
\end{align*}
左边是 $(T_1,\ldots,T_n)$ 在事件 $\{T_n<T\}$ 上的密度。任给
$0<t_1<\cdots<t_n$ 都可选择 $T>t_n$，故这就证明了整个正单纯形上的无条件密度
\eqref{eq:6-3-3-arrival-joint-density}。作三角变量替换
\[
 w_1=t_1,\qquad w_k=t_k-t_{k-1}\quad(k\geq2),
\]
雅可比行列式为 $1$，且 $t_n=\sum_{k=1}^nw_k$。故 $(W_1,\ldots,W_n)$ 的密度为
\[
 \lambda^ne^{-\lambda\sum_kw_k}\prod_{k=1}^n\1_{(0,\infty)}(w_k)
 =\prod_{k=1}^n\lambda e^{-\lambda w_k}\1_{(0,\infty)}(w_k).
\]
这对每个 $n$ 都成立，因而整列等待时间独立同分布为指数分布。特别地，$n=1$ 时退化为
$\Prob(T_1>t)=\Prob(N_t=0)=e^{-\lambda t}$。
\end{proof}

只有边缘正确仍可能构造出错误的过程分布。

\begin{example}[正确边缘、错误增量]\label{ex:6-3-3-4}
取一个 $U\sim\operatorname{Unif}(0,1)$，令 $\widetilde N_t=F_t^{-1}(U)$，其中 $F_t$ 是
$\operatorname{Poisson}(\lambda t)$ 的分布函数，分位数取左连续约定。每个
$\widetilde N_t$ 都有正确边缘，并且可取成随 $t$ 非减；可是全部时刻由同一个 $U$ 驱动。

例如取 $\lambda=1,s=1,t=2$。事件
$\{\widetilde N_1=0,\widetilde N_2=1\}$ 对应
\[
 e^{-2}<U\leq\min\{e^{-1},3e^{-2}\}=e^{-1},
\]
所以其概率为 $e^{-1}-e^{-2}$。真正的 Poisson 过程中，独立增量给
\[
 \Prob(N_1=0,N_2=1)=\Prob(N_1=0)\Prob(N_2-N_1=1)=e^{-2}.
\]
两数不同。这一比较直接说明一维 Poisson 公式不能替代独立增量。
\end{example}

中心化计数揭示线性漂移背后的鞅。

\begin{theorem}[Poisson 补偿鞅]\label{theorem:6-3-3-3}
设 $N$ 是速率 $\lambda$ 的 Poisson 过程，$M_t=N_t-\lambda t$。相对于原始自然滤过
$\mathcal F_t^N=\sigma(N_s:0\le s\le t)$ 或它的完备化，以下三个过程都是平方可积鞅：
\[
 M_t,\qquad M_t^2-\lambda t,\qquad M_t^2-N_t.
\]
\end{theorem}

\begin{proof}
固定 $0\leq s<t$，置 $Z=N_t-N_s$。它独立于 $\mathcal F_s$，且
$\E Z=\Var(Z)=\lambda(t-s)$。因此
\[
 \E[M_t-M_s\mid\mathcal F_s]=\E[Z-\lambda(t-s)]=0,
\]
证明 $M$ 为鞅。再令 $Y=Z-\lambda(t-s)$；$Y$ 独立于过去、均值为零、二阶矩为
$\lambda(t-s)$，所以
\[
 \E[M_t^2\mid\mathcal F_s]
 =\E[(M_s+Y)^2\mid\mathcal F_s]
 =M_s^2+\lambda(t-s).
\]
这给第二个鞅。最后
\[
 \E[M_t^2-N_t\mid\mathcal F_s]
 =M_s^2+\lambda(t-s)-N_s-\lambda(t-s)
 =M_s^2-N_s.
\]
Poisson 变量具有全部有限阶矩：由其概率质量函数，任意整数 $r\ge1$ 都有
\[
 \E N_t^r=e^{-\lambda t}\sum_{k=0}^\infty
 k^r\frac{(\lambda t)^k}{k!}<\infty,
\]
因为相邻项之比趋于零。特别地，前两个含 $M_t^2$ 的过程也属于 $L^2$，而上面的条件期望计算
均合法。把滤过完备化只加入零集及其子集；条件期望等式因此仍成立。本结论不需要预先断言原始
自然滤过已经右连续。
\end{proof}

$\lambda t$ 是由过去可以预测的平方波动累计量，而 $N_t$ 是实际跳跃平方和；二者之差仍是鞅。
这组公式将在 Brownian 二次变差出现后得到连续路径对照。

\begin{example}[复合 Poisson 过程]\label{ex:6-3-3-2}
令 $Y_1,Y_2,\ldots$ 为独立同分布的 $\R^d$ 值随机向量，共同分布为 $\nu$，并与 $N$ 独立。置
\[
 X_t=\sum_{k=1}^{N_t}Y_k,
\]
空和取零。若 $\varphi_Y(u)=\E e^{i\langle u,Y_1\rangle}$，则对 $u\in\R^d$，
\begin{align*}
 \E e^{i\langle u,X_t\rangle}
 &=\sum_{n=0}^\infty\Prob(N_t=n)\varphi_Y(u)^n\\
 &=e^{-\lambda t}\sum_{n=0}^\infty\frac{(\lambda t\varphi_Y(u))^n}{n!}
 =\exp\{\lambda t(\varphi_Y(u)-1)\}.
\end{align*}
这里 $e^{i\langle u,x\rangle}$ 只是正弦、余弦两项的合写，计算不使用复分析。
\end{example}

\begin{proposition}[复合 Poisson 过程的生成元]\label{proposition:6-3-3-compound-generator}
对从 $x$ 出发的复合 Poisson 过程，
\[
 P_tf(x)=\E f\left(x+\sum_{k=1}^{N_t}Y_k\right).
\]
对有界可测 $f$，其点态生成元为
\[
 Lf(x)=\lambda\int_{\R^d}\bigl(f(x+y)-f(x)\bigr)\nu(\dd y).
\]
若 $f\in C_0(\R^d)$，则 $P_tC_0\subset C_0$、
$\|P_tf-f\|_\infty\to0$，并且上述差商在 $C_0$ 范数中收敛。
\end{proposition}

\begin{proof}
当 $h\downarrow0$ 时，
\[
 \Prob(N_h=0)=1-\lambda h+O(h^2),\quad
 \Prob(N_h=1)=\lambda h+O(h^2),\quad
 \Prob(N_h\geq2)=O(h^2).
\]
按这三个事件分割，并用 $|f|\leq\|f\|_\infty$ 控制最后一项，得到一致于 $x$ 的展开
\[
 P_hf(x)=f(x)+\lambda h\int(f(x+y)-f(x))\nu(\dd y)+O(h^2)\|f\|_\infty.
\]
这给生成元公式。若 $f\in C_0$，则
$x\mapsto\int f(x+y)\nu(\dd y)$ 连续：对 $x_n\to x$ 用 DCT；它也在无穷远消失：先取紧集
$K$ 使 $\nu(K)>1-\varepsilon$，在 $K$ 上利用 $f(x+y)$ 的一致尾小性，在 $K^c$ 上用
$\|f\|_\infty\varepsilon$ 控制。Poisson 级数遂说明 $P_tf\in C_0$，而上面的统一展开给强连续性
与范数生成元结论。取 $f(x)=e^{iu\cdot x}$ 时，
\[
 Lf(x)=\lambda(\varphi_Y(u)-1)e^{iu\cdot x},
\]
与前一个特征函数计算相符。
\end{proof}

\begin{definition}[$Q$-矩阵]\label{def:6-3-3-3}
在可数状态空间 $E$ 上，数族 $(q_{ij})$ 称为保守 $Q$-矩阵，若
\[
 q_{ij}\geq0\ (i\ne j),\qquad
 q_i:=\sum_{j\ne i}q_{ij}<\infty,\qquad q_{ii}=-q_i.
\]
其形式生成元为
\[
 Qf(i)=\sum_jq_{ij}f(j)
 =\sum_{j\ne i}q_{ij}(f(j)-f(i)).
\]
在状态 $i$ 停留 $\operatorname{Exp}(q_i)$ 时间后，以概率 $q_{ij}/q_i$ 跳到 $j$；若 $q_i=0$，
约定 $i$ 为吸收态。
\end{definition}

\begin{example}[有限状态跳链]\label{ex:6-3-3-3}
若 $E$ 有限，取 $\Lambda>0$ 且 $\Lambda\geq\max_iq_i$，并置
\[
 R=I+\Lambda^{-1}Q.
\]
$R$ 是转移矩阵。取速率 $\Lambda$ 的 Poisson 时钟，并在每次时钟响起时按 $R$ 更新状态；
$R(i,i)=1-q_i/\Lambda$ 对应一次虚跳。条件于时钟次数，转移矩阵为
\[
 P_t=e^{-\Lambda t}\sum_{n=0}^\infty\frac{(\Lambda t)^n}{n!}R^n
 =e^{-\Lambda t}e^{\Lambda tR}
 =e^{\Lambda t(R-I)}=e^{tQ}.
\]
这称为一致化（uniformization）。Poisson 时钟在有限时间内只响有限次，故构造不爆炸。
\end{example}

同一论证适用于 $\sup_iq_i<\infty$ 的可数状态链，此时 $Q$ 是 $\ell^\infty(E)$ 上的有界算子。
仅有每个 $q_i<\infty$ 不够：纯出生链 $q_{i,i+1}=2^i$ 的连续等待时间期望和为
$\sum_i2^{-i}<\infty$，而非负总等待时间的期望有限迫使其总和 a.s. 有限；链可在有限时间内完成
无穷多次跳跃。这就是爆炸，必须另加条件排除。

\begin{proposition}[Kolmogorov 前后向方程]\label{proposition:6-3-3-4}
在有限状态或一致有界率情形，一致化定义的半群在矩阵范数或
$\ell^\infty$ 算子范数中可微，并满足
\[
 \frac{\dd}{\dd t}P_t=QP_t=P_tQ,\qquad P_0=I.
\]
\end{proposition}

\begin{proof}
$Q$ 有界，算子幂级数
$P_t=e^{tQ}=\sum_{n\geq0}t^nQ^n/n!$ 在每个有界时间区间上一致绝对收敛；导数级数也一致绝对
收敛。因此可逐项求导，得到
\[
 P_t'=\sum_{n=1}^\infty\frac{nt^{n-1}Q^n}{n!}
 =Q\sum_{m=0}^\infty\frac{t^mQ^m}{m!}=QP_t.
\]
$Q$ 与自己的幂交换，故同样等于 $P_tQ$。在概率语言中，左乘 $Q$ 是从起点作局部展开，右乘
$Q$ 是在终点作局部展开；二者由半群律协调。
\end{proof}

\begin{remark}[Poisson 与复合 Poisson 模型的应用]
若把 $N_t$ 解释为到达数，定理~\ref{theorem:6-3-3-2} 给出具有指数间隔的事件流；在排队模型中
它描述顾客到达。若 $Y_k\ge0$ 是第 $k$ 件事故的赔款，则例~\ref{ex:6-3-3-2} 的
$X_t=\sum_{k\le N_t}Y_k$ 是保险风险的累计损失。这两种解释都直接调用已证的独立平稳增量，
而不是仅凭单时刻 Poisson 边缘。

复合 Poisson 过程是纯跳、独立平稳增量过程的基本原型。
定理~\ref{theorem:6-3-3-3} 的补偿鞅 $N_t-\lambda t$ 把可预测强度与实际计数分开，为计数过程的随机积分提供基本的中心化噪声。
\end{remark}

\begin{remark}[\cadlagPath{}的空间边界]
Poisson 与一般跳链的自然路径有跳点，所以不能把它们的分布放到连续路径空间
$C([0,T])$ 上。适合的容器是 \cadlagPathSpace{} $D([0,T])$，其上通常配 Skorokhod 拓扑。
对显式构造的阶梯路径，当前的有限维分布计算已经足够。若要研究一般 $D([0,T])$-值随机元的弱收敛，
还需明确 Skorokhod 拓扑、其 Borel $\sigma$-代数以及相应的紧性判据。
\end{remark}

\begin{figure}[htbp]
\centering
\begin{tikzpicture}[x=1cm,y=1cm]
 \draw[book line,-{Stealth[length=2.4mm]}] (0,0)--(9.4,0) node[right,book label] {$t$};
 \foreach \x/\k in {1.1/1,2.7/2,4.8/3,7.6/4}{
   \draw[BookAccent,semithick] (\x,-0.12)--(\x,0.9);
   \node[book label,below] at (\x,-0.12) {$T_{\k}$};
 }
 \draw[decorate,decoration={brace,amplitude=4pt},BookWarm]
   (1.1,1.05)--(2.7,1.05) node[midway,above=4pt,book label] {$W_2$};
 \draw[BookWarning,semithick] (0,-0.65)--(1.1,-0.65)--(1.1,-0.35)--(2.7,-0.35)--
   (2.7,-0.05)--(4.8,-0.05)--(4.8,0.25)--(7.6,0.25)--(7.6,0.55)--(9.1,0.55);
 \node[book label] at (8.6,0.95) {$N_t$};
 \node[book warm node] at (4.7,-1.45) {局部率 $Q$\ $\xrightarrow{\ e^{tQ}\ }$\ 转移半群 $P_t$};
\end{tikzpicture}
\caption{等待时间决定到达点，到达点累计成单位跳阶梯；一致有界跳率经一致化生成半群。}
\label{fig:6-3-3-arrivals-and-generator}
\end{figure}

\begin{exercise}
\begin{enumerate}
 \item 从密度卷积重新推导式~\eqref{eq:6-3-3-gamma-density} 与 Poisson 计数分布。
 \item 证明条件于 $N_t=n$ 的到达点是 $n$ 个均匀点的次序统计量。
 \item 直接核验 $N_t-\lambda t$、$(N_t-\lambda t)^2-N_t$ 的鞅等式。
 \item 对纯出生链例证明无穷等待时间之和 a.s. 有限，并说明这一步为什么不能用于一致有界率链。
\end{enumerate}
\end{exercise}

Poisson 过程把时间积累为跳跃平方和。下一节转向连续路径噪声：Brownian motion 没有跳跃，
却仍有非零二次变差；这正是随机积分不能按经典有限变差积分处理的原因。
